\section{Global optimization in accuracy bits}\label{sec:accuracy}

Exact response coordinates for higher-degree local costs can generate
large algebraic fields. Approximation avoids assembling those fields.
The construction in this section approximates each scalar inverse by
rational polynomials, optimizes their compositions in a fixed number
of leader and dual variables, and recovers a rational leader with
controlled error. The requested accuracy enters through its number of
bits, rather than its reciprocal.

\subsection{Model, output, and approximation tools}\label{subsec:accuracy-model}

Let $X\subseteq[0,1]^r$ be a rational polytope, and let
$C\in\Q^{k\times N}$ and $U\in\Q^{s\times N}$ be independent of
the leader. The dimensions $r,k,s$ are fixed. The functions $b(x)$
and $\ell_i(x)$ are rational affine. Each rational polynomial marginal
$g_i$ is strictly increasing on $[0,1]$, with $g_i(0)=0$ and
$G_i=g_i(1)>0$. Its coefficients may have either sign. Write
$P=\max_i\deg g_i$ and $f_i(t)=\int_0^t g_i(u)\,du$.
An arbitrary strictly convex polynomial local cost is included by
subtracting its initial derivative from both its marginal and its
linear incentive. Additive local constants do not affect responses.

Put $R=\max\{1,\max_a\sum_i|U_{ai}|\}$ and
$W=[-R,R]^s$. The rational polynomial $\phi(x,w)$ is convex in
$w\in W$ for every follower-feasible leader. Convexity and strict
increase are promises; strict increase can alternatively be checked by
univariate sign methods. No positive lower bound on either local or
aggregate curvature is imposed. The follower minimizes
\begin{equation}\label{eq:accuracy-follower}
 F_x(z)=\sum_{i=1}^N[f_i(z_i)-\ell_i(x)z_i]+\phi(x,Uz),
 \qquad z\in[0,1]^N,\quad Cz\le b(x).
\end{equation}
Equality resources can be represented by pairs of inequalities, with
the resulting row count fixed. The follower-feasible leaders form the
explicit rational polytope $X_F$ of \cref{lem:resource-projection}.
At every $x\in X_F$, strict convexity gives a unique response $z(x)$.

Local marginals and the aggregate are densely encoded. Upper polynomials
may instead be supplied by explicit monomial lists in $(x,z)$; all their
terms and numerical total degrees count toward the complexity bound.
In particular no expansion in all $N$ follower variables is required.
Let $L_{\rm in}$ denote total rational input bit length and $D_{\rm all}$
the maximum numerical degree of all supplied polynomials.

\begin{theorem}[Accuracy-bit global optimization]\label{thm:accuracy-global}
For fixed $r,k,s$, follower-feasible leader existence for
\eqref{eq:accuracy-follower} is decidable in polynomial time. If
$X_F\ne\varnothing$, then for any rational polynomial upper objective
$H(x,z)$ and requested integer $B\ge1$, one can return a rational
$\widehat x\in X_F$ and a rational number $\widehat V$ such that
\[
 H(\widehat x,z(\widehat x))\le
       \min_{x\in X_F}H(x,z(x))+2^{-B},\qquad
 |\widehat V-\min_{x\in X_F}H(x,z(x))|\le2^{-B}.
\]
The bit complexity and rational output lengths are polynomial in
$(L_{\rm in},D_{\rm all},B)$. The minima are attained. No exact
common-field representation of the full response or optimum is promised.
Response-dependent upper constraints require the qualifications in
\cref{subsec:accuracy-upper}.
\end{theorem}

Scalar inverse resource allocation and logarithmic dependence on inverse
accuracy for the convex allocation objective are established
\cite{PatrikssonStromberg2015,HochbaumShanthikumar1990}.
Approximation of algebraic sums in fixed dimension also has an established
bit-model treatment \cite{Vigneron2014}. Here the objective being
optimized is an independent, possibly signed and nonconvex function of
the allocation response. The proof must control the entire response and
the feasibility of the returned leader, rather than only the follower's
objective gap. General polynomial bilevel relaxation methods provide a
complementary global-optimization approach \cite{JeyakumarEtAl2016}.

The clipped inverse $h_i$ of $g_i$ is zero for targets at most zero,
one for targets at least $G_i$, and $g_i^{-1}$ between them.
\Cref{thm:inverse-approximation} supplies rational polynomial branches
approximating $h_i$ to any tolerance $\eta$, with a complete uniform
error and polynomial-bit construction. Its proof in \cref{app:inverse}
handles real and complex critical values, interior derivative zeros,
rational analytic panels, and exact Taylor reversion.
\Cref{prop:positive-inverse} retains sharper bounds when coefficients
are nonnegative. We first explain how such approximations are used in
the diagonal and single-resource cases, before proving the full theorem.

\begin{lemma}[Rational recovery in a polytope]\label{lem:rational-rounding}
In fixed dimension $n$, let $Q\subseteq[0,1]^n$ be a nonempty rational
polytope and let $v\in Q$ have a polynomial-size common-field algebraic
encoding. For any rational $h>0$, a rational $\widehat v\in Q$ with
$\norm{\widehat v-v}_\infty\le h$ can be computed in polynomial bit
time in these encodings and $\max\{0,\log(1/h)\}$.
For a rational polynomial $S(v)=\sum_\alpha a_\alpha v^\alpha$,
\begin{equation}\label{eq:polynomial-lipschitz}
             L(S)=\max\{1,\sum_\alpha |a_\alpha|\,|\alpha|_1\}
\end{equation}
bounds its infinity-norm Lipschitz constant on the cube. Thus recovery
can preserve any polynomial-size finite list of polynomial values to
specified positive tolerances, while keeping membership in $Q$ exact.
\end{lemma}
\begin{proof}
Enumerate vertices by independent active-row systems and use the
convex-combination reduction in \cref{lem:core-recovery} to find an
affinely independent subset of at most $n+1$ vertices with
$v=\sum_{j=1}^t\lambda_jv_j$, $\lambda_j\ge0$, $\sum\lambda_j=1$.
Enumerating these subsets and checking barycentric coordinates takes
polynomial time in fixed dimension, in the existing algebraic field.
For $j<t$ replace $\lambda_j$ by $\lfloor M\lambda_j\rfloor/M$
and put the remaining mass on vertex $t$. All weights stay nonnegative;
the point stays in the same simplex and moves by at most $2n/M$.
Choose an integer $M\ge\max\{1,2n/h\}$. The floors require only
$O(\log M)$ exact algebraic comparisons each. Lower-dimensional and
singleton polytopes are included. Finally sum the absolute first partial
derivatives on the cube to obtain \eqref{eq:polynomial-lipschitz}, and
take $h$ no larger than the smallest tolerance divided by its bound.
\end{proof}

\subsection{Diagonal followers and dyadic powers}\label{subsec:accuracy-diagonal}

When $k=s=0$, the response is $z_i(x)=h_i(\ell_i(x))$.
For an affine upper objective $d_0+d\trans x+c\trans z(x)$, let
$\eps=2^{-B}$, $A=\sum_i|c_i|$ and
$\eta=\eps/[16\max\{1,A\}]$. If $A=0$, solve the leader LP.
Otherwise pull back every inverse breakpoint through $\ell_i$ and
construct the realizable affine sign cells in $X$, including zero signs.
There are polynomially many in fixed dimension. Take their closures:
weakening the strict affine tests gives the closure, since mixing any
point satisfying the weak tests with a point of the original relative
cell approaches it through that cell. Choose a valid branch on each
closed cell $Q$. Their polynomial upper objective $H_Q$ satisfies
\begin{equation}\label{eq:diagonal-uniform-error}
                     |H_Q(x)-H(x,z(x))|\le A\eta\le\eps/16.
\end{equation}
At a threshold either adjacent branch is valid on the boundary.

Minimize every $H_Q$ exactly on $Q$ by \cref{subsec:tools}, using
one additional copy of the $r$ variables to express that no point has a
smaller value. Compare the resulting algebraic values pairwise and
select the smallest. This requires at most products of two polynomial
field degrees, not a field containing all response radicals. If $v^*$
is the winning point, then $H_Q(v^*)\le\mathrm{OPT}+\eps/16$.
Use \cref{lem:rational-rounding} inside $Q$ so that its objective
increases by at most $\eps/4$. The true objective of the resulting
rational leader is at most $\mathrm{OPT}+3\eps/8$ by
\eqref{eq:diagonal-uniform-error}. Bisection of each true inverse at
that leader to weighted total error $\eps/4$ gives a rational optimum
estimate of error less than $\eps$. Rational vertex recovery is essential:
coordinatewise rounding of an algebraic minimizer can leave $X$.

For the important power marginal $g_i(t)=t^{p_i}$, one can use a simpler
fully explicit approximation. Put $m=\lceil\log_2(1/\eta)\rceil$,
$P=\max_i p_i$, $K_0=Pm$, and $q=m+1$. Below $2^{-K_0}$ use zero;
its root error is at most $2^{-m}$. On
$2^{-(j+1)}\le t\le2^{-j}$, $j=0,\ldots,K_0-1$, put
$b_j=3\cdot2^{-(j+2)}$ and $u=t/b_j-1\in[-1/3,1/3]$.
For $\alpha=1/p_i$, the rational binomial coefficients have absolute
value at most one, by their consecutive ratio. Thus
\[
 T_q(u)=\sum_{n=0}^q\binom\alpha n u^n,\quad
 |T_q(u)|\le3/2,\quad
 |(1+u)^\alpha-T_q(u)|\le3^{-q}/2\le\eta/4.
\]
Compute a rational $\sigma_j$ within $\eta/4$ of $b_j^{1/p_i}$ by
bisection and use $\sigma_jT_q(t/b_j-1)$. Its error is at most
$3\eta/8+\eta/4<\eta$. There are $Pm$ dyadic bands and degree
$O(m)$. Bisection compares rational $p_i$th powers of
$O(m)$-bit arguments against $O(Pm)$-bit dyadic targets. Binomial
coefficient lengths are $O(q(\log P+\log q))$, and expansion adds
$O(qK_0)$ bits from scaling. These are polynomial in numerical $P$
and accuracy bits. In fixed leader dimension the substituted polynomial
has at most $\binom{r+q}{r}$ monomials. This proves the diagonal
accuracy-bit guarantee directly, without a critical-value calculation.

\begin{example}[Sparse exponents can force long rational output]
\label{ex:sparse-power-output}
Let even $p\ge2$ be encoded in binary. Two independent unit-box
followers with marginals $z^{p}$ and $z^{p/2}$ and incentive $x\in[0,1]$
have responses $x^{1/p}$ and $x^{2/p}$. Their difference has minimum
$-1/4$ at $x=2^{-p}$, since it is $(x^{1/p}-1/2)^2-1/4$.
Suboptimality at most $1/64$ implies
$3/8\le x^{1/p}\le5/8$, hence $0<x\le(5/8)^p$.
For reduced rational $x=a/b>0$, this requires $b\ge(8/5)^p$,
so explicit rational output takes $\Omega(p)$ bits, although the
sparse input has $O(\log p)$ bits and accuracy is constant. The
example uses two different powers. It is an output obstruction, not a
hardness claim for the same-power subclass.
\end{example}

\subsection{One equality and an exact signed balance identity}
\label{subsec:accuracy-one-resource}

For $s=0$ with one resource equality $\sum_iw_iz_i=b(x)$, where
$w_i\in\Q$ may be negative or zero, there is a sharper argument than
the general residual estimate below. Define
\[
 b_-=\sum_i\min\{w_i,0\},\quad b_+=\sum_i\max\{w_i,0\},\quad
 W_1=\sum_i|w_i|.
\]
The box image is exactly $[b_-,b_+]$: its endpoint vertices and their
segment realize that interval. Thus $X_F=X\cap\{b_-\le b(x)\le b_+\}$.
If $W_1=0$, this is simply $b(x)=0$ and the diagonal construction
applies. Otherwise put $w_{\min}=\min_{w_i\ne0}|w_i|$.

Bound each affine incentive on the leader cube by rational
$\ell_i^-,\ell_i^+$. Among all $i$ with $w_i\ne0$, take the minimum
$\lambda_-$ and maximum $\lambda_+$ of
\[
 \ell_i^-/w_i,\quad\ell_i^+/w_i,\quad
 (\ell_i^--G_i)/w_i,\quad(\ell_i^+-G_i)/w_i.
\]
They have polynomial bit length and $\lambda_-<\lambda_+$.
The box-Lagrangian response and its balance are
\[
 q_i(x,\lambda)=h_i(\ell_i(x)-\lambda w_i),\qquad
 B_x(\lambda)=\sum_iw_iq_i(x,\lambda).
\]
The continuous nonincreasing function $B_x$ equals $b_+$ at
$\lambda_-$ and $b_-$ at $\lambda_+$. Hence every feasible leader
has a balancing multiplier in that interval. Its box minimizer is the
resource optimum by Lagrangian comparison against any balanced point.
This includes endpoint resources and requires no constraint qualification.
All balancing multipliers give the same response by strict convexity.

For a balancing $\lambda^*$, every nonzero term
$w_i(q_i(x,\lambda)-q_i(x,\lambda^*))$ has the same weak sign.
Zero-weight responses do not change. Therefore
\begin{equation}\label{eq:signed-balance-identity}
 \sum_i|w_i|\,|q_i(x,\lambda)-z_i(x)|
                         =|B_x(\lambda)-b(x)|.
\end{equation}
In particular, for $c_{\max}=\max_i|c_i|$,
\[
 |c\trans(q-z(x))|\le(c_{\max}/w_{\min})|B_x(\lambda)-b(x)|.
\]
This identity uses monotonicity and the single dual parameter, including
signed weights. There is no corresponding no-cancellation identity for
several resource rows.

Here is the complete affine-objective accuracy ledger for this case.
If $c=0$ solve the LP over $X_F$. Otherwise let
$A=\sum_i|c_i|$, $C_1=c_{\max}W_1/w_{\min}$ and
$\eta=\eps/[32(1+A+C_1)]$. Normalize
$\lambda=\lambda_-+(\lambda_+-\lambda_-)\theta$,
$0\le\theta\le1$, and arrange the affine inverse breakpoints in
$X_F\times[0,1]$. On each closed rational cell $Q$ let $p_i$ be its
valid approximants, and put
\[
 H_Q=d_0+d\trans x+c\trans p,\qquad T_Q=\sum_iw_ip_i-b(x).
\]
Minimize $H_Q$ subject to $v=(x,\theta)\in Q$ and
$|T_Q(v)|\le W_1\eta$. A true optimal leader and a balancing multiplier
have such a lift, so the winning algebraic value is at most
$\mathrm{OPT}+A\eta$. Optimize in $2(r+1)$ variables and recover a
rational point in $Q$ within distance
\[
 h\le\min\{1,\eps/(8L(H_Q)),W_1\eta/L(T_Q)\}.
\]
The recovered polynomial residual is at most $2W_1\eta$ and the
true box response residual at most $3W_1\eta$. Consequently
\[
 |H(\widehat x,z(\widehat x))-H_Q(\widehat v)|
      \le(A+3C_1)\eta,
\]
and its suboptimality is at most
$\eps/8+(2A+3C_1)\eta\le7\eps/32$. The rational polynomial value
at $\widehat v$ lies between
$\mathrm{OPT}-(A+3C_1)\eta$ and
$\mathrm{OPT}+A\eta+\eps/8$, hence estimates the optimum within
$5\eps/32$. The leader is exactly in $X_F$ even though its auxiliary
response need not be balanced. Small weights affect logarithmic precision
only. Compact balancing multipliers also prove response continuity by
subsequence passage and uniqueness, and therefore attainment.

\subsection{Resource and aggregate residuals control the response}
\label{subsec:accuracy-certificate}

We now return to \eqref{eq:accuracy-follower}, assuming $N\ge1$.
Construct $X_F$ and $K$ by \cref{lem:resource-projection,lem:resource-constants}.
Choose a rational $\overline G\ge1$ bounding
$\norm{\nabla_zF_x(z)}_\infty$ on the leader and follower cubes.
Absolute coefficient sums after normalizing $w=R(2t-\mathbf1)$
supply such a bound with polynomial encoding. Put
\[
 \begin{gathered}
 L_F=N\overline G,\qquad \Lambda=\max\{1,K\overline G\},\\
 S=\max\{1,\max_j\sum_i|C_{ji}|\},\qquad
 A_U=\max\{1,\max_a\sum_i|U_{ai}|\}=R.
 \end{gathered}
\]
Empty row maxima are zero. The constant $L_F$ bounds the follower
objective's infinity-norm Lipschitz constant. Let $L_\phi\ge0$ bound
the infinity-norm Lipschitz constant of $\nabla_w\phi$ on $W$,
uniformly on the leader cube; the maximum row sum of absolute Hessian
coefficient bounds suffices. Define
\[
 D_\phi=L_\phi\sum_i\norm{U_i}_1,\qquad
                        \mu=\frac{\min_iG_i}{4(4P)^P}.
\]
These are effective rational constants of polynomial bit length. With
positive marginal coefficients one may use the larger $\mu$ from
\cref{lem:polynomial-bregman}.

For $x\in X_F$, $w\in W$ and $\lambda\in[0,\Lambda]^k$, define
the frozen-gradient box response
\begin{equation}\label{eq:accuracy-frozen-response}
 q_i=h_i\bigl(\ell_i(x)-U_i\trans\nabla_w\phi(x,w)
                                      -C_i\trans\lambda\bigr).
\end{equation}
This minimizes a strictly convex separable box objective with aggregate
gradient frozen at $w$.

\begin{lemma}[Response certificate]\label{lem:accuracy-certificate}
If the true frozen response satisfies
\begin{equation}\label{eq:accuracy-residuals}
 Cq-b(x)\le\delta\mathbf1,\quad
 |\lambda\trans(Cq-b(x))|\le\zeta,\quad
                       \norm{Uq-w}_\infty\le\rho,
\end{equation}
where $\delta,\zeta,\rho\ge0$, then
\begin{equation}\label{eq:accuracy-response-bound}
 \norm{q-z(x)}_\infty
 \le K\delta+\left(\frac{L_FK\delta+\zeta+D_\phi\rho}{\mu}
                                                      \right)^{1/(P+1)}.
\end{equation}
An exact response has a certificate with all three residuals zero and
$\lambda\in[0,\Lambda]^k$. The response $z(x)$ is continuous on $X_F$.
\end{lemma}
\begin{proof}
Put $a=\nabla_w\phi(x,w)$ and $a_q=\nabla_w\phi(x,Uq)$.
Testing the frozen box variational inequality at $z^*=z(x)$ and
using convexity of the actual follower objective gives
\[
 F_x(z^*)\ge F_x(q)-\lambda\trans C(z^*-q)
                            +(a_q-a)\trans U(z^*-q)
 \ge F_x(q)+\lambda\trans(Cq-b(x))-D_\phi\rho.
\]
Here $Cz^*\le b(x)$, $\lambda\ge0$, and
$\norm{U(z^*-q)}_1\le\sum_i\norm{U_i}_1$. Repair $q$ to feasible
$y$ within $K\delta$ by \cref{lem:resource-constants}. Then
\[
 0\le F_x(y)-F_x(z^*)\le L_FK\delta+\zeta+D_\phi\rho.
\]
First-order optimality at $z^*$ and \cref{lem:polynomial-bregman}
bound $\mu\norm{y-z^*}_\infty^{P+1}$ by this gap. Add the repair
distance to prove \eqref{eq:accuracy-response-bound}.

At an exact optimum the bounded multipliers of
\cref{lem:resource-constants} and $w=Uz^*$ give
\eqref{eq:accuracy-frozen-response} with zero residuals. Conversely
zero residuals imply the full convex follower KKT conditions and identify
its unique response. Along a sequence in $X_F$, the aggregate and
multiplier boxes are compact; every convergent subsequence of such
certificates passes to a zero-residual certificate by continuity of
the clipped inverses. Uniqueness identifies every response limit.
This proves continuity, including at degenerate feasible faces.
\end{proof}

The aggregate error has a useful sharper exponent. Retaining
$\norm{z^*-q}_\infty$ in the preceding dot-product bound and writing
$t=\norm{y-z^*}_\infty$ yields
\[
 \mu t^{P+1}\le L_FK\delta+\zeta+D_\phi\rho(K\delta+t).
\]
If $t$ exceeded both terms in the following maximum, each right-hand
term would be less than half the left side. Therefore
\begin{equation}\label{eq:accuracy-sharp-mismatch}
 \norm{q-z^*}_\infty\le K\delta+
 \max\left\{
 \left(\frac{2(L_FK\delta+\zeta+D_\phi\rho K\delta)}\mu\right)^{1/(P+1)},
 \left(\frac{2D_\phi\rho}\mu\right)^{1/P}\right\}.
\end{equation}
When $\delta=\zeta=0$, take $y=q$ and divide by the response distance
if it is positive to obtain the sharper
$\norm{q-z^*}_\infty\le(D_\phi\rho/\mu)^{1/P}$.
The simpler bound \eqref{eq:accuracy-response-bound} suffices below.
For $s=0$ it reduces to the multiple-resource residual/complementarity
bound, with no scalar balance identity presumed.

\subsection{Polynomial response approximations and rational recovery}
\label{subsec:accuracy-graph}

Normalize $w=R(2t-\mathbf1)$, $\lambda=\Lambda\theta$, so the
compressed variable $v=(x,t,\theta)$ lies in the rational polytope
$Q_0=X_F\times[0,1]^{s+k}$. Its dimension $r+s+k$ is fixed.
The inverse arguments
\[
 a_i(v)=\ell_i(x)-U_i\trans\nabla_w\phi(x,R(2t-\mathbf1))
                                           -\Lambda C_i\trans\theta
\]
are rational polynomials. For tolerance $0<\eta\le1/4$, pull back
all rational inverse breakpoints by these arguments and compute their
realizable sign conditions on $Q_0$ using \cref{subsec:tools}.
At each sign sample choose a valid branch for every inverse. Define
$Q$ by $Q_0$ and the \emph{closed branch-interval conditions} for that
chosen branch vector. These compact sets cover $Q_0$, and on each,
\begin{equation}\label{eq:accuracy-valid-branches}
                     |p_i(v)-q_i(v)|\le\eta.
\end{equation}
There are polynomially many sets, since there are polynomially many
realizable signs in fixed dimension. They may overlap or be disconnected.
The closed validity set may enlarge its originating sign condition,
but every added point retains all branch error bounds. This avoids
both a Cartesian enumeration of branch choices and an unjustified
closure operation on nonlinear cells. When $s=0$, the arguments are
affine and the sets can be the closed rational arrangement cells.

On each $Q$ define the polynomial residuals
\[
 E=Cp-b(x),\qquad J=\lambda\trans E,\qquad T=w-Up,
\]
and the compact candidate set
\begin{equation}\label{eq:accuracy-candidate}
 \mathcal S_Q=\{v\in Q:E_j\le S\eta,\quad
        |J|\le k\Lambda S\eta,\quad |T_a|\le A_U\eta\}.
\end{equation}
Empty row families are omitted. Every true response has a lift in some
candidate set, with $p$ within $\eta$ of that response: use its true
aggregate and a bounded multiplier. The polynomial residual allowances
follow from \eqref{eq:accuracy-valid-branches}. Every candidate in turn
has true residual allowances $2S\eta$, $2k\Lambda S\eta$ and
$2A_U\eta$. All polynomials have polynomial degree and coefficient
length after expansion in the fixed compressed dimension.

\begin{lemma}[Recovery across nonlinear branch boundaries]
\label{lem:accuracy-recovery}
Given $0<\tau\le1/4$, choose
\begin{equation}\label{eq:accuracy-general-tolerance}
 \eta\le\min\left\{\frac14,\tau,\frac{\tau}{8KS},
 \frac{\mu(\tau/2)^{P+1}}
              {1+4L_FKS+5k\Lambda S+4D_\phi A_U}\right\}.
\end{equation}
For any algebraically encoded $v^*\in\mathcal S_Q$ and any rational
$h_0>0$, one can compute a rational $\widehat v\in Q_0$ with
$\norm{\widehat v-v^*}_\infty\le h_0$ such that
\begin{equation}\label{eq:accuracy-recovery-interface}
 \norm{q(\widehat v)-q(v^*)}_\infty\le\eta,\qquad
 \norm{q(\widehat v)-z(\widehat x)}_\infty\le\tau,\qquad
 \norm{p(v^*)-z(\widehat x)}_\infty\le2\eta+\tau.
\end{equation}
All additional precision and bit cost are polynomial in the supplied
encodings, in $\log(1/\tau)$, and in the positive part of
$\log(1/h_0)$.
Membership in the nonlinear branch set $Q$ is not promised.
\end{lemma}
\begin{proof}
Set
\[
 \alpha=\min_i\frac{G_i}{2}\left(\frac{\eta}{2P}\right)^P.
\]
By \cref{lem:inverse-modulus}, an argument change less than $\alpha$
changes every clipped inverse by less than $\eta$: for actual degree
$d\le P$, $(\eta/(2d))^d\ge(\eta/(2P))^P$.
Let $M_a\ge1$ bound the Lipschitz constants of all $a_i$ on the
normalized cube. Let $B_b\ge1$ bound the affine Lipschitz constant of
$b$ and let $\overline E=S+\sup_{x\in[0,1]^r}\norm{b(x)}_\infty$,
using an absolute coefficient upper bound for the supremum. Choose
positive rational $h$ strictly smaller than
\begin{equation}\label{eq:accuracy-rounding-distance}
 \min\left\{h_0,\frac{\alpha}{2M_a},
 \frac{S\eta}{B_b+1},\frac{A_U\eta}{2R+1},
                                      \frac{S\eta}{\overline E+1}\right\}.
\end{equation}
Use \cref{lem:rational-rounding} in $Q_0$. Then true clipped responses
change by at most $\eta$, right-hand sides by at most $S\eta$,
aggregates by at most $A_U\eta$, and multipliers by at most
$\Lambda S\eta/(\overline E+1)$.

At $v^*$ the true residual allowances were twice those in
\eqref{eq:accuracy-candidate}. At $\widehat v$ they are therefore
\begin{equation}\label{eq:accuracy-transferred-residuals}
       \delta=4S\eta,\qquad \zeta=5k\Lambda S\eta,
                                      \rho=4A_U\eta.
\end{equation}
For complementarity write, with the \emph{true} residuals $e^*,\widehat e$,
\[
 \widehat\lambda\trans\widehat e-(\lambda^*)\trans e^*
 =\widehat\lambda\trans(\widehat e-e^*)
                         +(\widehat\lambda-\lambda^*)\trans e^*.
\]
The first term has magnitude at most $2k\Lambda S\eta$ and the
second at most $k\Lambda S\eta$, because
$\norm{e^*}_\infty\le\overline E$. This absolute residual bound
is needed for inactive constraints with large negative slack.
\Cref{lem:accuracy-certificate} and
\eqref{eq:accuracy-general-tolerance} bound the repair and growth terms
by $\tau/2$ each, proving \eqref{eq:accuracy-recovery-interface}.
Every displayed constant and precision has polynomial bit length;
raising a tolerance to $P+1$ costs polynomially many bits in numerical
$P$. The proof compares true continuous clipped responses across the
boundary. It never evaluates the old branch polynomial at a point where
that branch might be invalid.
\end{proof}

For resource-only followers a sharper, entirely polyhedral recovery
is available. Choose
\begin{equation}\label{eq:accuracy-resource-tolerance}
 \eta\le\min\left\{\frac14,\tau,\frac{\tau}{6SK},
 \frac{\mu(\tau/2)^{P+1}}{3S(1+L_FK+k\Lambda)}\right\}.
\end{equation}
Round within its rational branch cell, allowing each $E_j$ to change
by at most $S\eta$ and $J$ by at most $k\Lambda S\eta$.
\Cref{lem:rational-rounding} also permits any finite list of additional
polynomial rounding tolerances. The true residuals at the rounded
point are then $3S\eta$ and $3k\Lambda S\eta$, so
\cref{lem:accuracy-certificate} gives $\norm{q-z(\widehat x)}_\infty\le\tau$.
Thus the valid rounded polynomial response obeys
$\norm{p(\widehat v)-z(\widehat x)}_\infty\le\eta+\tau$.
For an affine objective, set $\tau=\eps/[16(1+\norm{c}_1)]$ and
allow objective rounding error $\eps/8$. A true optimum has polynomial
lift value at most $\mathrm{OPT}+\norm{c}_1\eta$, while the final
true-to-polynomial error is at most
$\norm{c}_1(\eta+\tau)\le\eps/8$.
The final suboptimality is at most $5\eps/16$; the rational rounded
value lies between $\mathrm{OPT}-\eps/8$ and
$\mathrm{OPT}+3\eps/16$. This is the complete fixed-resource
ledger, including equalities and vanishing curvature. It uses repair
and polynomial growth, not \eqref{eq:signed-balance-identity}.

\subsection{Polynomial upper functions and the global theorem}
\label{subsec:accuracy-polynomial-upper}

For an explicitly listed upper polynomial
$G(x,z)=\sum_{\alpha,\beta}a_{\alpha\beta}x^\alpha z^\beta$, put
\begin{equation}\label{eq:upper-polynomial-lipschitz}
 \begin{split}
 L_x(G)&=\sum_{\alpha,\beta}|a_{\alpha\beta}|\,|\alpha|_1
                                                        2^{|\beta|_1},\\
 L_z(G)&=\sum_{\alpha,\beta}|a_{\alpha\beta}|\,|\beta|_1
                                                        2^{|\beta|_1}.
 \end{split}
\end{equation}
These bound the separate infinity-norm Lipschitz constants on
$[0,1]^r\times[-1,2]^N$, by summing absolute partial derivatives.
They may vanish when their variable group is absent. The enlarged
response box includes approximants with $\eta\le1/4$.
Their encodings are polynomial in the explicit input and numerical
upper degree, including the factor $2^{|\beta|_1}$.

On each candidate substitute $z=p(v)$. If the inverse branches have
degree at most $d_p$, the resulting polynomial has degree at most
$D_{\rm up}\max\{1,d_p\}$ in the fixed compressed variables.
Successive polynomial multiplication expands every listed monomial in
polynomial time: only polynomially many monomials exist in fixed
dimension at that degree, and coefficient growth is polynomial in
numerical degree and input coefficient lengths. The number of upper
polynomials and terms can grow without increasing the optimization
dimension.

For the recovery in \cref{lem:accuracy-recovery}, the segment between
$p(v^*)$ and $z(\widehat x)$ lies in $[-1,2]^N$, so
\begin{equation}\label{eq:upper-recovery-error}
 |G(\widehat x,z(\widehat x))-G(x^*,p(v^*))|
 \le L_x(G)h_0+L_z(G)(2\eta+\tau).
\end{equation}
At a true response lift the corresponding error is at most $L_z(G)\eta$.

\begin{proof}[Proof of \cref{thm:accuracy-global}]
Projection decides emptiness. Continuity from
\cref{lem:accuracy-certificate} and compactness of $X_F$ ensure
attainment. Let $\eps=2^{-B}$,
$A_H=\max\{1,L_x(H),L_z(H)\}$, and choose
$\tau\le\eps/(32A_H)$ and $h_0\le\eps/(32A_H)$, with
$\eta$ as in \eqref{eq:accuracy-general-tolerance}.
Minimize the substituted upper polynomial on every candidate set
\eqref{eq:accuracy-candidate} and compare their exact minima.
Each minimizing-point formula uses two copies of $r+s+k$ variables,
so \cref{subsec:tools} gives polynomial-size algebraic minima
and samples. At a true optimizer the lift exists, so the winning
value $V$ satisfies $V\le\mathrm{OPT}+A_H\eta\le\mathrm{OPT}+\eps/32$.
Recover the winner by \cref{lem:accuracy-recovery}.
Equation \eqref{eq:upper-recovery-error} gives
$|H(\widehat x,z(\widehat x))-V|\le4A_H\eps/(32A_H)=\eps/8$.
Exact membership $\widehat x\in X_F$ also gives
$V\ge\mathrm{OPT}-\eps/8$. Approximate $V$ rationally to error
$\eps/16$ to obtain the claimed value estimate and the leader guarantee.
Only polynomial-degree algebraic outputs in fixed dimension are compared;
no exact sum of the true inverse values is evaluated.

If $N=0$, the follower constraints reduce to $0\le b(x)$ and there
is no response. Minimize the polynomial upper objective on this rational
polytope in fixed dimension and apply \cref{lem:rational-rounding}.
This is an LP only when the upper objective is affine. All counts,
degrees, coefficient encodings and requested precisions above are
polynomial in $(L_{\rm in},D_{\rm all},B)$.
\end{proof}

For example, a convex quartic aggregate may couple many units through
one total output while local second derivatives vanish at an optimum.
The objective $-x_1\sum_i z_i+d(x)$ includes tariff revenue and a
polynomial intervention cost. Neither follower convexity nor the theorem
requires the induced revenue to be concave. These are structural examples;
the coarse constants and fixed-dimensional elimination do not establish
an efficient implemented solver.

\subsection{Upper constraints: approximation and exact feasibility}
\label{subsec:accuracy-upper}

Let $G_j(x,z)\le0$, $j=1,\ldots,m$, be explicitly listed rational
polynomial upper inequalities, with $m$ unrestricted. Put
$\mathcal H(x)=H(x,z(x))$, $\mathcal G_j(x)=G_j(x,z(x))$ and
\begin{equation}\label{eq:tightened-value}
 V(t)=\min\{\mathcal H(x):x\in X_F,\ \mathcal G_j(x)\le-t
                  \text{ for all }j\},\qquad t\ge0,
\end{equation}
with $V(t)=+\infty$ for an empty set. By continuity and compactness
each finite minimum is attained. Leader-only affine rows should be
placed directly in $X$ and preserved exactly. Response-dependent
upper equalities may be represented by opposite inequalities, but in
general cannot satisfy the positive tightening assumptions below.

Margin conversion from response error to upper feasibility is standard,
including in near-optimal robust bilevel models
\cite{BesanconAnjosBrotcorne2024}. The value of perturbed inequalities
and convex interpolation toward a strict point are also established
\cite[Section 5.6]{BoydVandenberghe2004}. Here they are combined with
the proved global accuracy-bit construction and its rational recovery.

\begin{theorem}[Outer and inner guarantees]\label{thm:accuracy-upper}
For positive rational objective and violation tolerances $\eps,\delta$,
there are algorithms with bit complexity polynomial in the input,
numerical degrees, and the encodings of these tolerances, for fixed
$r,k,s$, with the following properties.

The \emph{outer} algorithm either certifies original infeasibility, or
returns rational $\widehat x\in X_F$ with
$\mathcal G_j(\widehat x)\le\delta$ for every $j$.
If $V(0)$ is finite it returns a point and
$\mathcal H(\widehat x)\le V(0)+\eps$.
An outer point alone does not certify original feasibility.

The \emph{inner} algorithm either reports an empty inner surrogate, or
returns rational $\widehat x\in X_F$ satisfying all original upper
inequalities exactly. If $V(\delta)$ is finite it returns a point with
$V(0)\le\mathcal H(\widehat x)\le V(\delta)+\eps$.
An empty inner surrogate implies $V(\delta)=+\infty$, but does not
certify original infeasibility.
\end{theorem}
\begin{proof}
Let $A$ be the maximum of one and all $L_x,L_z$ constants for the
objective and upper rows, and choose
\[
 \kappa=\min\{1/4,\eps/(16A),\delta/(16A)\},\qquad
 \tau\le\kappa,\quad h_0\le\kappa,
\]
with $\eta\le\kappa$ satisfying
\eqref{eq:accuracy-general-tolerance}. On each candidate set write
$F_j(v)=G_j(x,p(v))$ and $H_Q(v)=H(x,p(v))$.
A true lift changes each criterion by at most $L_z(G_j)\eta$.
At any candidate and its recovery, \eqref{eq:upper-recovery-error}
is at most $4A\kappa\le\min\{\eps,\delta\}/4$.

For the outer algorithm impose $F_j\le L_z(G_j)\eta$ on every
candidate set and minimize $H_Q$ exactly. Every true feasible leader
has an admissible lift. Hence emptiness certifies original infeasibility,
and, when $V(0)$ is finite, the winning value $V_{\rm out}$ is at
most $V(0)+L_z(H)\eta$. Recover its algebraic minimizer using
\cref{lem:accuracy-recovery}. The true violation is at most
$\delta/16+\delta/4$ and objective loss at most
$\eps/16+\eps/4$, which prove the stated bounds.

For the inner algorithm impose $F_j\le-\delta/2$.
Every leader feasible for $V(\delta)$ has an admissible lift, since
$L_z(G_j)\eta\le\delta/16$. Any recovered candidate satisfies
$\mathcal G_j(\widehat x)\le-\delta/4<0$. If $V(\delta)$ is
finite, its lift bounds the winning polynomial value by
$V(\delta)+L_z(H)\eta$, and recovery gives the stated objective
bound. All optimization still uses a fixed number of variables even
though $m$ grows. In the resource-only case one may instead use the
rational-cell recovery and its additional polynomial rounding tolerances;
the same conclusions follow from its smaller response-error bounds.

When $N=0$, optimize the upper polynomials directly over
$X\cap\{0\le b(x)\}$ with the same outer/inner margins and rational
polytope recovery. This is not in general an LP. Likewise, a
response-independent objective does not justify skipping response
approximation when any upper constraint still depends on it.
\end{proof}

The two algorithms give a posteriori global certificates without any
regularity promise. Let rational $v_{\rm out},v_{\rm in}$ approximate
the exact winning surrogate values within $\omega_{\rm out},\omega_{\rm in}$.
For the inner recovery put
$e_{\rm in}=L_x(H)h_0+L_z(H)(2\eta+\tau)$, using its own tolerances.
If both procedures return, then
\begin{equation}\label{eq:accuracy-posterior}
 \underbrace{v_{\rm out}-\omega_{\rm out}-L_z(H)\eta_{\rm out}}_{\mathrm{Lower}}
 \le V(0)\le\mathcal H(x_{\rm in})\le
 \underbrace{v_{\rm in}+\omega_{\rm in}+e_{\rm in}}_{\mathrm{Upper}}.
\end{equation}
The inner point supplies the needed original feasibility. An outer
objective alone is not an absolute-error estimate of $V(0)$ and can
lie below it. As all tolerances tend to zero, the outer lower bounds
converge to finite $V(0)$: otherwise compactness, vanishing violation,
and continuity yield a feasible limit with smaller objective. Inner
upper bounds also converge if $V(t)\to V(0)$ as $t\downarrow0$.
The computed finite gap is always a valid certificate when both
bounds exist; a quantitative stopping guarantee needs more.

\begin{corollary}[A supplied tightening modulus]
\label{cor:accuracy-tightening}
Suppose $V(0)<+\infty$ and supplied rationals $D_0,\sigma>0$ and
integer $\nu\ge1$ satisfy
\begin{equation}\label{eq:tightening-modulus}
 V(t)\le V(0)+D_0(t/\sigma)^{1/\nu},\qquad0<t\le\sigma.
\end{equation}
Then an exactly upper-feasible rational leader within additive $\eps$
of $V(0)$, and a rational absolute-error optimum estimate, are computable
in polynomial time in the input, numerical degrees, numerical $\nu$,
and accuracy bits. The modulus parameters enter through their encodings;
the algorithm need not verify the promise.
\end{corollary}
\begin{proof}
Use inner tightening
$t=\sigma\min\{1/2,(\eps/(2D_0))^\nu\}$ and objective tolerance
$\eps/2$. The modulus ensures nonemptiness and tightening loss at
most $\eps/2$. Its rational encoding is polynomial in numerical
$\nu$ and the supplied bit lengths. Repeat with smaller constant
fractions of the requested error and approximate the inner achieved
objective using its surrogate/recovery bound to obtain an absolute-error
estimate of $V(0)$.
\end{proof}

\begin{corollary}[Convex reduced constraints with a strict anchor]
\label{cor:accuracy-anchor}
Suppose each $\mathcal G_j$ is convex on $X_F$, and a supplied point
$x_s\in X_F$ has $\mathcal G_j(x_s)\le-\sigma$ for all $j$, with
rational $\sigma>0$. For $N\ge1$, \cref{cor:accuracy-tightening} applies with
\[
 \nu=P,\qquad D_0=\max\{1,L_x(H)+L_z(H)C_z\},
\]
where the effective constant $C_z$ is given in
\cref{prop:accuracy-response-modulus}. One can alternatively use its
simpler modulus with exponent $\nu=P+1$ and the corresponding constant.
For $N=0$, use $\nu=1$ and $D_0=\max\{1,L_x(H)\}$ by direct
Lipschitz interpolation. Convexity of the reduced constraints and the anchor margin are additional
promises, not consequences of follower convexity.
\end{corollary}
\begin{proof}
At an original optimizer $x^*$, set
$x_t=(1-t/\sigma)x^*+(t/\sigma)x_s$, for $0\le t\le\sigma$.
It lies in $X_F$, has all reduced constraints at most $-t$, and moves
by at most $t/\sigma$. Equation \eqref{eq:response-modulus-sharp}
and \eqref{eq:upper-polynomial-lipschitz} bound its objective change
by $D_0(t/\sigma)^{1/P}$. This proves the modulus.
The algorithm uses the supplied margin; it does not silently solve an
exact algebraic-sum comparison to certify the anchor.
\end{proof}

For a concrete independent-follower subclass, assume the increasing
marginals are convex, $g_i(0)=0$, and the affine incentives are
nonnegative throughout $X$. Then
$z_i(x)=\min\{1,g_i^{-1}(\ell_i(x))\}$ is concave: the increasing
inverse is concave and its constant extension above saturation preserves
concavity. Thus $a_{j0}+a_j\trans x-\sum_iw_{ji}z_i(x)$ is convex
for $w_{ji}\ge0$, including minimum weighted production constraints.
The incentive sign matters: clipping a negative incentive to zero can
destroy concavity. This example does not assert concavity under shared
resources or aggregate coupling.

\begin{corollary}[Reserve controls with uniform headroom]
\label{cor:accuracy-reserve}
Write the leader as $(u,s)\in X_u\times S_r$, with a fixed-dimensional
rational reserve polytope $S_r\subseteq[0,1]^p$. Suppose the complete
follower objective and constraints depend only on $u$, and at least one
follower-feasible $u$ exists. For each feasible
$u$, assume every $G_j(u,s,z(u))$ is convex in $s$, and there is a
supplied $s_{\rm safe}\in S_r$ with uniform margin
$G_j(u,s_{\rm safe},z(u))\le-\sigma$ for all $u,j$.
Then \cref{cor:accuracy-tightening} applies with $\nu=1$ and
$D_0=\max\{1,L_s(H)\}$, where the polynomial coefficient bound for
reserve derivatives supplies $L_s(H)$.
\end{corollary}
\begin{proof}
Keep the incentive component $u^*$ of an original optimizer and
interpolate its reserve toward $s_{\rm safe}$. The follower response
is unchanged. Convexity in the reserve gives margin $t$ at interpolation
fraction $t/\sigma$, and the objective changes by at most
$L_s(H)t/\sigma$. The safe reserve also supplies original feasibility
whenever the follower-feasible $u$ set is nonempty.
\end{proof}

For affine rows $a_{j0}+a_j\trans u+c_j\trans z-R_js\le0$,
uniform headroom can be certified conservatively by rational LP over
the explicit polytope $X_{u,F}\times[0,1]^N$ with $s=s_{\rm safe}$.
The matrix $R$ need not have a prescribed sign. For an objective
$h(u)+w\trans s$, use $D_0=\max\{1,\norm{w}_1\}$. Reserves are
actual model controls, such as external procurement, not artificial
slacks inserted into an existing hard constraint. Their dimension counts
toward the fixed leader dimension. The aggregate must also be independent
of them. No convexity of the reduced constraints in $u$ is needed.

\subsection{Why the feasibility qualifications are necessary}
\label{subsec:accuracy-limits}

\begin{example}[A strict point does not ensure inner convergence]
\label{ex:accuracy-isolated-optimum}
For $x,z\in[0,1]$, let
\[
 F_x(z)=z^2/2+z^4/4-z^3/6-xz,\qquad
 g(z)=z+z^2(z-1/2).
\]
The marginal derivative is
$g'(z)=3(z-1/6)^2+11/12>0$, so even uniform strong curvature holds.
Impose the affine upper row $z(x)-x\le0$ and minimize $x$.
Since $g(0)=0$, $g(1)=3/2$, and $x\in[0,1]$, the response is its
unclipped inverse. Monotonicity gives
\[
 z(x)\le x\ \Longleftrightarrow\ x\le g(x)
                  \ \Longleftrightarrow\ x^2(x-1/2)\ge0.
\]
Thus the feasible set is $\{0\}\cup[1/2,1]$ with optimum zero and
strict feasible point one. Every positively tightened feasible leader
has $x>1/2$. By continuity and strict feasibility for $x>1/2$,
$V(t)\to1/2$ as $t\downarrow0$, so the original optimum is not
approached. A strict point alone cannot justify an inner convergence
or tightening-modulus claim.
\end{example}

\begin{example}[No rational exactly feasible leader]
\label{ex:accuracy-irrational-feasibility}
The follower cost $z^3/3-xz$ on the unit box gives $z(x)=\sqrt x$.
The affine upper equality $z-x=1/8$ permits precisely
\[
                         x=3/8\pm\sqrt2/4.
\]
Both roots lie in $[0,1]$ and satisfy the unsquared equation, but both
are irrational. Thus feasible rational leader output is impossible
without qualifications, even with one leader and one follower.
\end{example}

Finally, constant incentives $a_i/M^2$, cubic local costs and upper
weights $M$, with integer $M\ge\max_i a_i>0$, give the fixed upper
value $\sum_i\sqrt{a_i}$. One exact upper inequality therefore
contains sum-of-square-roots comparison. The accuracy algorithms do not
resolve this exact arithmetic problem: outer constraints allow prescribed
violation and exact inner feasibility uses stated margins. Numerical
degree dependence is also essential by \cref{ex:sparse-power-output}.
Neither these arithmetic limits nor the rational-output examples are
claims that a convex follower cannot itself be approximated efficiently.
